% !TeX root = ../教材题库_全部条目_转写工作稿.tex
% Source: XJ-B1-C05-HW-P162-Y1089
% 第5章 / 5.1 任意角和弧度制 / 教材 p.162 / PDF 第169页
% OCR未校对：公式与文字均可能有误；未完成数学TeX转写；禁止直接发给学生
% 原题及相邻原解截图（核对用）：../题目原图/XJ-B1-C05-HW-P162-Y1089-01.webp
% 本文件只在转写完成后作为题干使用；OCR 条目尚不是可靠题干。
\begin{bankitem}{习题5.1 第11题}{第5章 / 5.1 任意角和弧度制 / 教材 p.162 / PDF 第169页}{XJ-B1-C05-HW-P162-Y1089}{OCR未校对：公式与文字均可能有误；未完成数学TeX转写；禁止直接发给学生}
\begin{ocrdraft}
[教材p.162 / PDF第169页]
\par
II.已知圆上按顺时针排列的5个点A，B,C，D，E将该圆分为五段弧，其
\par
长度之比，：．CD：：万乁一1：2：3：4：5，求五边形ABCDE的各个内
\par
角的弧度数．
\end{ocrdraft}
\end{bankitem}
